\section{Teoretický základ}\label{sec:teoreticky-zaklad}

\subsection{Binárne obrázky a ich reprezentácia}\label{subsec:binarne-obrazky-a-ich-reprezentacia}

\subsubsection{Definícia binárnych obrázkov}

\hspace{0.5cm}
Binárny obrázok je sekvencia dat v binárnej podobe, kde pixely nadobúdajú iba dve možné hodnoty 0 alebo 1.
Zvyčajne čiernou farbou je interpretované pozadie a je reprezentované nulami a popredie alebo objekt zaujmú je biele,
reprezentované jednotkami.

\subsubsection{Problém dekompozície na obdĺžniky}

\hspace{0.5cm}
Problém dekompozície na obdĺžniky je postup, pri ktorom sa popredie obrázka (pixely s hodnotou 1) rozloží na
množinu navzájom neprekrývajúcich sa ortogonálnych obdĺžnikov tak, aby ich zjednotenie presne prekrývalo všetky
pixely popredia a žiadne iné.

Formálne: Nech $B$ je binárny objekt (binárnym objektom rozumieme množinu všetkých pixelov binárneho obrázka,
ktorých hodnota je rovná~1).
Rozložíme ho na $K \geq 1$ blokov $B_1, B_2, \ldots, B_K$ takých,
že $B_i \cap B_j = \emptyset$ pre každé $i \neq j$ a $B = \bigcup_{k=1}^{K} B_k$. \parencite{SUK20124279}

Existuje viacero metód, ako sa dá rozložiť binárny objekt na obdĺžniky, ich algoritmy sa trocha líšia a sú zamerané nie len
na samotný rozklad, ale aj dosiahnuť rozklad s čo menším počtom blokov alebo dosiahnuť najlepšiu kompresiu obrázka v rozumnom čase.
Tento problém je NP-ťažký, čo bolo dokázané viacerými autormi \parencite{hanauer2024covering,masek1978np,najman2014covering}

Pre všeobecné binárne obrazy neexistuje známy polynomiálny algoritmus na nájdenie optimálneho riešenia.
Napriek tomu je možné získať optimálne riešenia tohto problému použitím lineárneho celočíselného programovania (ILP),
ktoré poskytuje presné výsledky.


\subsection{Genetické algoritmy}
\label{subsec:geneticke-algoritmy}

\subsubsection{Princípy evolučných algoritmov}

\hspace{0.5cm}
Genetický algoritmus, (ďalej GA) je vyhľadávacia heuristika inšpirovaná prirodzeným výberom a evolúciou,
ktorá sa používa na nájdenie optimálnych riešení zložitých problémov.

\subsubsection{Reprezentácia chromozómov}

\subsubsection{Fitness funkcia}

\subsubsection{Genetické operátory}

\subsubsection{Stratégie selekcie}



